\subsection{为什么稳态流动仍然需要外力}
稳态表示固定位置的流速不随时间变化，不表示经过这里的每一份流体都没有加速度。图~\ref{ch13:fig:jet-plate} 中，虚线区域内的水总量可以不变，但新进入的水不断丢失法向动量，这就需要平板持续施力。
在短时间 $\dd t$ 内，流入质量为 $\dot m\dd t$，流入动量为 $\dot m\vect v_{\rm in}\dd t$；流出端同理。稳态时控制体内的动量不积累，所以外力冲量必须补足出入动量之差。除以 $\dd t$ 得到上面的动量流方程。多出口时必须先按矢量相加，不能直接相加速率。

对管内流动，在入口截面，外部流体把所研究流体推入控制体；在出口截面，外部流体阻碍它流出。两个压力方向均指向控制体内部。若求管道与环境之间的净作用力，使用表压时还须把管外大气压的抵消计入，不能只把绝对压强随意置零。

\begin{origexamplebox}{射流推动匀速后退的平板}\label{ch13:ex:teaching-moving-plate}
沿水平 $+x$ 方向的细水射流速度为 $v$、截面积为 $A$、密度为 $\rho$。一块足够大的平板沿同方向以恒速 $u$ 后退，$0\le u<v$。水撞板后沿板面流走，相对平板没有法向速度。忽略重力及空气作用，求水对板的力和传给板的功率。
\end{origexamplebox}
\begin{xesolution}
以匀速平板为参考系，它仍是惯性系。对应图~\ref{ch13:fig:jet-plate} 的控制体随板平移，入射相对速度为 $v-u$，故每秒真正撞板的质量为
\[
 \dot m_{\rm hit}=\rho A(v-u).
\]
这不是固定截面的 $\rho Av$：一部分水还在追赶平板。水相对板的法向速度从 $v-u$ 变成零，每单位质量丢失法向动量 $v-u$。由反作用关系，板受到沿 $+x$ 的力
\[
 F=\dot m_{\rm hit}(v-u)=\rho A(v-u)^2.
\]
也可在地面系看同一份水，撞前法向速度为 $v$，撞后为 $u$，其动量变化给出同一个结果。力对平板的功率必须使用地面系板速：
\[
 P=Fu=\rho Au(v-u)^2.
\]
在 $0<u<v$ 内，$\dd P/\dd u=\rho A(v-u)(v-3u)$，故最大功率出现在 $u=v/3$，大小为 $4\rho Av^3/27$。$u=0$ 时有力却不向静止板传机械功；$u\to v$ 时几乎没有水追上平板，力和功率均趋零。这也说明力最大与功率最大不是同一个条件。
\end{xesolution}
